\documentclass[12pt, letterpaper]{article}
\usepackage[margin=0.75in, top=0.75in, bottom=0.75in]{geometry}
\usepackage{booktabs}
\usepackage{longtable}
\usepackage{array}
\usepackage{xcolor}
\usepackage{colortbl}
\usepackage{hyperref}
\usepackage{helvet}
\renewcommand{\familydefault}{\sfdefault}

% Color definitions
\definecolor{headerbg}{RGB}{31, 73, 125}
\definecolor{headertext}{RGB}{255, 255, 255}
\definecolor{row1}{RGB}{235, 241, 250}
\definecolor{row2}{RGB}{255, 255, 255}
\definecolor{domainA}{RGB}{198, 217, 241}
\definecolor{domainB}{RGB}{217, 234, 211}
\definecolor{domainC}{RGB}{255, 230, 199}
\definecolor{domainD}{RGB}{234, 209, 220}

\hypersetup{
    colorlinks=true,
    linkcolor=blue,
    urlcolor=blue
}

\begin{document}

\begin{center}
    {\LARGE \textbf{4th Grade Math --- Full Year Curriculum}}\\[4pt]
    {\large 25 Lessons across 10 Units}
\end{center}

\vspace{6pt}

\small
\begin{longtable}{
    >{\centering\arraybackslash}p{0.4cm}   % \#
    >{\raggedright\arraybackslash}p{3.0cm}  % Domain
    >{\raggedright\arraybackslash}p{2.8cm}  % Unit
    >{\raggedright\arraybackslash}p{3.2cm}  % Lesson Title
    >{\raggedright\arraybackslash}p{5.2cm}  % Description
    >{\centering\arraybackslash}p{2.0cm}    % Standard
}

% Header
\rowcolor{headerbg}
\textcolor{headertext}{\textbf{\#}} &
\textcolor{headertext}{\textbf{Domain}} &
\textcolor{headertext}{\textbf{Unit}} &
\textcolor{headertext}{\textbf{Lesson Title}} &
\textcolor{headertext}{\textbf{Description}} &
\textcolor{headertext}{\textbf{Standard}} \\
\endfirsthead

\rowcolor{headerbg}
\textcolor{headertext}{\textbf{\#}} &
\textcolor{headertext}{\textbf{Domain}} &
\textcolor{headertext}{\textbf{Unit}} &
\textcolor{headertext}{\textbf{Lesson Title}} &
\textcolor{headertext}{\textbf{Description}} &
\textcolor{headertext}{\textbf{Standard}} \\
\endhead

\multicolumn{6}{r}{\small\itshape Continued on next page\ldots} \\
\endfoot

\bottomrule
\endlastfoot

\toprule

% ── DOMAIN 1: Operations & Algebraic Thinking ──────────────────────────────────
\multicolumn{6}{l}{\cellcolor{domainA}\textbf{Domain: Operations \& Algebraic Thinking}} \\
\midrule

\rowcolor{row1}
1 & Operations \& Algebraic Thinking & Unit 1: Multiplicative Comparison & Comparison Word Problems &
Interpret a multiplication equation as a comparison (e.g., $35 = 5 \times 7$ means 35 is 5 times as many as 7). Solve word problems using multiplicative comparisons. &
4.OA.A.1 / 4.OA.A.2 \\

\rowcolor{row2}
2 & Operations \& Algebraic Thinking & Unit 1: Multiplicative Comparison & Multi-Step Word Problems &
Solve multi-step word problems involving whole numbers and all four operations. Represent word problems using equations with a variable. &
4.OA.A.3 \\

\rowcolor{row1}
3 & Operations \& Algebraic Thinking & Unit 2: Factors, Multiples & Patterns & Factors \& Multiples &
Find all factor pairs for a whole number in the range 1--100. Recognize multiples of numbers and define prime vs.\ composite numbers. &
4.OA.B.4 \\

\rowcolor{row2}
4 & Operations \& Algebraic Thinking & Unit 2: Factors, Multiples & Patterns & Generating Patterns &
Generate a number or shape pattern that follows a given rule (e.g., add 3). Analyze patterns to identify informal features. &
4.OA.C.5 \\

\rowcolor{row1}
5 & Operations \& Algebraic Thinking & Unit 2: Factors, Multiples & Patterns & Equations \& Unknowns &
Solve equations with variables for unknown numbers in multiplication, division, addition, and subtraction word contexts. &
4.OA.A.3 (applied) \\

\midrule
% ── DOMAIN 2: Place Value & Multi-Digit Arithmetic ────────────────────────────
\multicolumn{6}{l}{\cellcolor{domainB}\textbf{Domain: Place Value \& Multi-Digit Arithmetic}} \\
\midrule

\rowcolor{row1}
6 & Place Value \& Multi-Digit Arithmetic & Unit 3: Multi-Digit Place Value & Place Value Structure &
Recognize that in a multi-digit whole number, a digit in one place represents ten times what it represents in the place to its right. Read and write numbers. &
4.NBT.A.1 / 4.NBT.A.2 \\

\rowcolor{row2}
7 & Place Value \& Multi-Digit Arithmetic & Unit 3: Multi-Digit Place Value & Rounding \& Comparing Numbers &
Use place value understanding to round multi-digit whole numbers to any place. Compare large numbers using $<$, $>$, and $=$. &
4.NBT.A.2 / 4.NBT.A.3 \\

\rowcolor{row1}
8 & Place Value \& Multi-Digit Arithmetic & Unit 4: Computation & Multi-Digit Addition \& Subtraction &
Fluently add and subtract multi-digit whole numbers using the standard algorithm. Solve real-world word problems. &
4.NBT.B.4 \\

\rowcolor{row2}
9 & Place Value \& Multi-Digit Arithmetic & Unit 4: Computation & Multi-Digit Multiplication &
Multiply a whole number of up to four digits by a one-digit whole number, and multiply two two-digit numbers. Illustrate with area models. &
4.NBT.B.5 \\

\rowcolor{row1}
10 & Place Value \& Multi-Digit Arithmetic & Unit 4: Computation & Multi-Digit Division &
Find whole-number quotients and remainders with up to four-digit dividends and one-digit divisors, using equations and rectangular arrays. &
4.NBT.B.6 \\

\midrule
% ── DOMAIN 3: Fractions & Decimals ──────────────────────────────────────────
\multicolumn{6}{l}{\cellcolor{domainC}\textbf{Domain: Fractions \& Decimals}} \\
\midrule

\rowcolor{row1}
11 & Fractions \& Decimals & Unit 5: Fraction Equivalence & Equivalent Fractions &
Explain fraction equivalence by using visual fraction models. Generate equivalent fractions by multiplying or dividing numerators and denominators. &
4.NF.A.1 \\

\rowcolor{row2}
12 & Fractions \& Decimals & Unit 5: Fraction Equivalence & Comparing Fractions &
Compare two fractions with different numerators and different denominators by creating common denominators or comparing to benchmarks. &
4.NF.A.2 \\

\rowcolor{row1}
13 & Fractions \& Decimals & Unit 6: Fraction Arithmetic & Adding \& Subtracting Fractions &
Understand addition and subtraction of fractions as joining and separating parts referring to the same whole. Add/subtract like fractions and mixed numbers. &
4.NF.B.3a-b \\

\rowcolor{row2}
14 & Fractions \& Decimals & Unit 6: Fraction Arithmetic & Decomposing & Word Problems &
Decompose a fraction into a sum of fractions with the same denominator in more than one way. Solve word problems with fraction addition/subtraction. &
4.NF.B.3c-d \\

\rowcolor{row1}
15 & Fractions \& Decimals & Unit 6: Fraction Arithmetic & Multiplying Fractions by Whole Numbers &
Understand a fraction as a multiple of a unit fraction (e.g., $5/4 = 5 \times 1/4$). Multiply a fraction by a whole number in word problems. &
4.NF.B.4 \\

\rowcolor{row2}
16 & Fractions \& Decimals & Unit 7: Decimal Fractions & Decimal Fractions &
Express a fraction with denominator 10 as an equivalent fraction with denominator 100, and use this technique to add two fractions. &
4.NF.C.5 \\

\rowcolor{row1}
17 & Fractions \& Decimals & Unit 7: Decimal Fractions & Decimal Notation \& Comparison &
Use decimal notation for fractions with denominators 10 or 100. Compare two decimals to hundredths by reasoning about their size ($<$, $>$, $=$). &
4.NF.C.6 / 4.NF.C.7 \\

\midrule
% ── DOMAIN 4: Measurement & Data ───────────────────────────────────────────────
\multicolumn{6}{l}{\cellcolor{domainD}\textbf{Domain: Measurement \& Data}} \\
\midrule

\rowcolor{row1}
18 & Measurement \& Data & Unit 8: Measurement Conversions & Measurement Size Relationships &
Know relative sizes of measurement units (km, m, cm; kg, g; lb, oz; l, ml; hr, min, sec). Record measurement equivalents in a two-column table. &
4.MD.A.1 \\

\rowcolor{row2}
19 & Measurement \& Data & Unit 8: Measurement Conversions & Measurement Word Problems &
Use the four operations to solve word problems involving distances, intervals of time, liquid volumes, masses of objects, and money. &
4.MD.A.2 \\

\rowcolor{row1}
20 & Measurement \& Data & Unit 8: Measurement Conversions & Area \& Perimeter of Rectangles &
Apply the area and perimeter formulas for rectangles in real-world and mathematical problems (e.g., finding width given area). &
4.MD.A.3 \\

\rowcolor{row2}
21 & Measurement \& Data & Unit 9: Representing Data \& Angles & Line Plots with Fractions &
Make a line plot to display a data set of measurements in fractions of a unit ($1/2$, $1/4$, $1/8$). Solve addition/subtraction fraction problems. &
4.MD.B.4 \\

\rowcolor{row1}
22 & Measurement \& Data & Unit 9: Representing Data \& Angles & Understanding Angle Measurement &
Recognize angles as geometric shapes formed wherever two rays share a common endpoint. Understand concept of angle measurement in degrees. &
4.MD.C.5 \\

\rowcolor{row2}
23 & Measurement \& Data & Unit 9: Representing Data \& Angles & Measuring \& Adding Angles &
Measure angles in degrees using a protractor. Sketch angles of specified measure. Solve addition/subtraction problems finding unknown angles. &
4.MD.C.6 / 4.MD.C.7 \\

\midrule
% ── DOMAIN 5: Geometry & Financial Literacy ───────────────────────────────
\multicolumn{6}{l}{\cellcolor{domainA}\textbf{Domain: Geometry \& Financial Literacy}} \\
\midrule

\rowcolor{row1}
24 & Geometry \& Financial Literacy & Unit 10: Shape Lines & Angles & Lines, Angles \& Shape Symmetry &
Draw points, lines, line segments, rays, angles (right, acute, obtuse), and perpendicular and parallel lines. Identify line-symmetric shapes. &
4.G.A.1 / 4.G.A.3 \\

\rowcolor{row2}
25 & Geometry \& Financial Literacy & Unit 10: Shape Lines & Angles & Classifying Shapes \& Financial Literacy &
Classify two-dimensional figures based on the presence or absence of parallel or perpendicular lines. Solve money problems including tax and change. &
4.G.A.2 / 4.MD.A.2 (applied) \\

\bottomrule
\end{longtable}

\end{document}
